% !TEX root = ../main.tex
\section{Diseño del registro Estudiante}

El enunciado de la investigación pide administrar trece datos por
estudiante: número de carné, identificación, nombre, apellidos, fecha de
nacimiento, sexo, carrera, nivel académico, correo electrónico, número
telefónico, dirección, cantidad de créditos aprobados y promedio general,
además del estado del estudiante. El enunciado no propone ninguna forma de
organizar estos datos, solo los enumera como requisitos que el sistema
debe guardar.

El primer intento del equipo fue juntar los trece campos en un único
registro. Al revisarlo con cuidado apareció un problema: esos trece campos
no cambian con la misma frecuencia ni se usan juntos en las mismas
operaciones. El nombre y la fecha de nacimiento de un estudiante casi
nunca cambian una vez matriculado. El promedio y los créditos aprobados se
actualizan varias veces durante la carrera del estudiante, cada período
lectivo. El correo, el teléfono y la dirección se consultan juntos
normalmente solo cuando hace falta contactar al estudiante, y casi nunca
al mismo tiempo que se actualiza el promedio.

Juntar estos tres grupos en un solo registro significa que actualizar el
promedio de un estudiante, algo que ocurre constantemente, obliga también
a leer y reescribir su nombre, su fecha de nacimiento, su correo y su
dirección, aunque ninguno de esos datos haya cambiado.

Por esa razón se decidió dividir el registro en tres structs
independientes, cada uno guardado en su propio archivo binario
(\texttt{identidad.dat}, \texttt{academico.dat}, \texttt{contacto.dat}).
El criterio de separación es la frecuencia y el patrón de uso de cada
grupo de campos, no una división arbitraria. Esta técnica se conoce como
particionamiento vertical y es la misma que aplican los motores de bases
de datos relacionales cuando separan columnas de una entidad lógica en
tablas físicas distintas por razones de rendimiento.

Los tres structs se mantienen alineados por posición: el registro en la
posición $i$ de \texttt{identidad.dat} corresponde al mismo estudiante que
el registro en la posición $i$ de \texttt{academico.dat} y de
\texttt{contacto.dat}. Gracias a esto, un único índice (carné apuntando a
posición) basta para ubicar al estudiante en cualquiera de los tres
archivos, sin necesidad de tres índices repetidos.

\subsection{IdentidadEstudiante}
Datos que identifican al estudiante como persona: carné, identificación,
nombre, apellidos, fecha de nacimiento y sexo.

\begin{lstlisting}[style=pseudocodigo]
typedef struct {
    int carne;
    char identificacion[13];
    TipoIdentificacion tipoIdentificacion;
    char nombre[21];
    char apellido1[21];
    char apellido2[21];
    Fecha fechaNacimiento;
    int sexo;
} IdentidadEstudiante;
\end{lstlisting}

\subsection{AcademicoEstudiante}
Datos de la situación académica actual del estudiante, que cambian cada
período lectivo.

\begin{lstlisting}[style=pseudocodigo]
typedef struct {
    int carne;
    int carrera;
    int nivelAcademico;
    int creditosAprobados;
    float promedio;
    int estado;
} AcademicoEstudiante;
\end{lstlisting}

\subsection{ContactoEstudiante}
Datos usados para comunicarse con el estudiante, que cambian con poca
frecuencia.

\begin{lstlisting}[style=pseudocodigo]
typedef struct {
    int carne;
    char correo[41];
    char telefono[9];
    char direccion[81];
} ContactoEstudiante;
\end{lstlisting}
